% =====================================================================================
% SCOPE §13 project report -- TFIM with logical (encoded) vs physical qubits, aimed at
% Quantinuum H2.  Compiles with pdfLaTeX (Overleaf) using standard packages only.
%
% Figures: the five PNGs are produced by `scripts/plot_q3.py` into the repo's `figures/`
% folder (gitignored).  The macro \maybefig below looks for them in ../../figures/ (when
% compiled inside the repo) and then in ./figures/ (Overleaf: upload the PNGs into a
% folder called `figures` next to this file).  If neither exists, a framed placeholder
% is printed instead, so the document always compiles.
%
% Status-tag convention (project-wide): [established] checked source; [derived] proved
% here and tested; [measured] output of this repo's code; [inferred] reasoned, not
% tested; [assumed] a declared modelling choice; [unverified] not checked.
% Every noisy number is a CLASSICAL PREDICTION UNDER AN ASSUMED NOISE MODEL, not an
% H2 result.  No H2 experiment was run by this project.
% =====================================================================================
\documentclass[11pt,a4paper]{article}

\usepackage[T1]{fontenc}
\usepackage[utf8]{inputenc}
\usepackage{lmodern}
\usepackage[margin=2.3cm]{geometry}
\usepackage{amsmath,amssymb}
\usepackage{graphicx}
\usepackage{booktabs}
\usepackage{longtable}
\usepackage{array}
\usepackage{tabularx}
\usepackage{xcolor}
\usepackage{enumitem}
\usepackage[hidelinks]{hyperref}
\usepackage{xurl}% lets long file paths break anywhere

\setlist{itemsep=2pt,topsep=4pt}
\setlength{\LTcapwidth}{\linewidth}
\renewcommand{\arraystretch}{1.15}
\newcolumntype{L}{>{\raggedright\arraybackslash}X}
\newcolumntype{P}[1]{>{\raggedright\arraybackslash}p{#1}}

% ---- status tags -------------------------------------------------------------------
\newcommand{\tg}[1]{\textsf{\small[#1]}}
\newcommand{\est}{\tg{established}}
\newcommand{\der}{\tg{derived}}
\newcommand{\mea}{\tg{measured}}
\newcommand{\infr}{\tg{inferred}}
\newcommand{\asm}{\tg{assumed}}
\newcommand{\unv}{\tg{unverified}}
\newcommand{\pend}{\tg{pending}}
\newcommand{\cp}{\tg{classical prediction, assumed noise -- NOT H2}}

\newcommand{\code}[1]{\nolinkurl{#1}}
\newcommand{\Hcal}{\mathcal{H}}

% ---- figures that may be missing ------------------------------------------------------
\newcommand{\maybefig}[2]{%
  \IfFileExists{../../figures/#1}{\includegraphics[width=#2]{../../figures/#1}}{%
  \IfFileExists{figures/#1}{\includegraphics[width=#2]{figures/#1}}{%
  \fbox{\parbox{0.8\linewidth}{\centering\small Figure \code{#1} not found.\\
  Generate it with \code{python -I scripts/plot_q3.py} (writes \code{figures/}),\\
  or upload it to a \code{figures/} folder next to this file.}}}}}

\newenvironment{statusbox}{\begin{center}\begin{minipage}{0.94\linewidth}\hrule\vspace{4pt}\small}%
{\vspace{4pt}\hrule\end{minipage}\end{center}}

\title{Logical versus physical qubits for the transverse-field Ising model on Quantinuum H2:\\
\large classical groundwork, frustration, thermalization and QEC choice\\[2pt]
\normalsize Project report (SCOPE \S13)}
\author{TFIM logical-qubit project (K.~Martinez)\\
\small Prepared by the project's Claude Code overseer from the repository's reviewed deliverables}
\date{8 October 2026}

\begin{document}
\maketitle

\begin{statusbox}
\textbf{What this report is, and is not.}
Every number below comes from this repository's code and data files, or from a cited source; every load-bearing
claim in the body carries a status tag (the abstract summarizes tagged results). \textbf{No experiment on Quantinuum H2 (hardware or emulator) has been run by this project.}
Every curve or number obtained under noise is a \emph{classical prediction under an assumed noise model},
labelled \cp{} or ``assumed noise''. The only H2 measurements quoted are the \emph{published} physical-qubit data
of Haghshenas \emph{et al.}~\cite{haghshenas}, used as the reference; they are always labelled as such.
The project's headline logical-vs-physical comparison is \textbf{hardware-gated} and has not been performed.
\end{statusbox}

\begin{abstract}
We report the classical groundwork for an article whose thesis is to compute the \emph{same} physical
observables of the transverse-field Ising model (TFIM) with \emph{logical} (error-detected) qubits on
Quantinuum H2 and compare them head-to-head with the released physical-qubit data of
arXiv:2503.20870~\cite{haghshenas}. The contribution is framed as logical-vs-physical methodology,
frustration physics and QEC choice, not as beyond-classical dynamics: local Trotterized 2D Ising observables
at these sizes are classically reachable~\cite{begusic}.
(i)~\emph{QEC choice (Q4).} Four codes were built and run end to end (Iceberg $[[k{+}2,k,2]]$, a concatenated
Iceberg $[[16,4,4]]$, the 5-qubit $[[5,1,3]]$ and Steane $[[7,1,3]]$ codes). In every one the undetectable residual
sits at the arbitrary-angle logical rotation gadgets, not at the code distance, so the distance-3 and distance-4
codes buy only $\approx1.2\times$ on that residual at 3--4$\times$ lower acceptance at $p_2=10^{-3}$ (1.4--1.6$\times$ at
the H2 prior), under assumed noise. Iceberg is recommended.
(ii)~\emph{Does encoding help (Q1)?} Under an assumed noise model, this depends on the ZZ fraction $\eta$ of the
native two-qubit gate error; on the $2\times2$ test lattice (the most pessimistic case) the crossover is
$\eta\approx0.25$, rising to $\approx0.44$--0.55 on large square and triangular tori. At a conservative prior read from
the reference experiment's benchmarking ($\eta\approx0.05$) encoding is \emph{predicted} to help (Iceberg residual
$0.60\times$ bare per step on $2\times2$), conditional on randomized compiling. This is a classical prediction under
assumed noise, not a measurement.
(iii)~\emph{Frustration (Q2).} On the matched $4\times4$ pair, the triangular nearest-neighbour connected
correlation is 3--5$\times$ weaker than on the bipartite square lattice; an older square ``reference'' on a
$3\times4$ torus turned out to be itself frustrated and is corrected.
(iv)~\emph{Thermalization (Q3).} A sector-restricted Monte Carlo baseline is built; ideal Trotter curves relax
to the finite-cluster diagonal ensemble rather than the infinite-temperature sector value; the working step
$dt=0.25$ sits 1.3--2$\times$ below a measured Trotter threshold for $h\le0.3$; under the assumed noise,
Gauss-law post-selection is the cheap effective mitigation, and noise can mimic thermalization.
(v)~\emph{Emergent Gauss law.} As a resource for error mitigation it is mostly \emph{no}: a narrow,
conditional \emph{yes} for post-selection against bit-flip-type noise, floored at $O((h/J)^2)$.
A findings-vs-questions table (\S\ref{sec:table}) maps each conclusion to its evidence.
\end{abstract}

\tableofcontents

% =====================================================================================
\section{Introduction and thesis}\label{sec:intro}

\paragraph{The gap.} Haghshenas \emph{et al.}\ (arXiv:2503.20870, Nature 653, 56 (2026))~\cite{haghshenas}
simulated TFIM dynamics on up to 56 physical trapped-ion qubits of Quantinuum H2 and released the raw
observable-level data on Zenodo. Separate work runs logical-qubit simulations but does not compute
these physical observables from logical data. The project's thesis is to close that gap with a
\emph{controlled} comparison: the same Hamiltonian, the same observables and conventions, and encoding as
the single varied axis --- between this project's own encoded and bare arms, run on one lattice at matched
shots. The released experiment's physical-qubit data serves as an \emph{external reference} rather than as
the bare arm: Iceberg needs $k+4$ physical qubits, so the reference's $7\times8$ lattice would need 60 and H2
has 56, and lattice size, calibration and mitigation would all differ as well (see \S\ref{sec:q1} and
\S\ref{sec:hw}).
A cross-platform comparison was considered and rejected because it changes Hamiltonian, code, observable
and platform at once.

\paragraph{Related work, to be distinguished (not baselines).}
Froland \emph{et al.}~\cite{froland} report an encoded $[[4,2,2]]$ Iceberg vs unencoded comparison of Ising
local observables on IBM superconducting hardware \tg{established, abstract}. Quantinuum~\cite{helios} shows encoded logical
qubits beyond break-even on Helios, including an XY-model magnetism simulation \tg{established, abstract}. Neither is
matched to 2503.20870, uses the TFIM on H2, or treats the frustrated triangular lattice. Concatenated Iceberg
itself is not new (it is used in~\cite{helios}). The residual novelty is therefore the combination: H2 trapped ions,
TFIM matched to 2503.20870's released data (with Iceberg and concatenated $d=4$ encodings), and the frustrated
triangular lattice and its dimer observables; the last is the least pre-empted.

\paragraph{Honest framing (K4).} Local observables of Trotterized 2D Ising dynamics at 50--130 qubits have been
reproduced classically with sparse-Pauli and tensor-network methods~\cite{begusic} \est{}, corroborated by
belief-propagation tensor networks~\cite{tindall} \tg{fetched}. The contribution is
\emph{not} quantum advantage from the dynamics; it is logical-vs-physical methodology on real hardware
observables, frustration physics, and QEC methodology.

\paragraph{Priority and questions.} SCOPE \S2 fixes four questions in priority order Q1 (logical vs physical) $>$
Q2 (frustration) $>$ Q3 (thermalization), plus Q4 (QEC choice) and an extra skeptical thread on the emergent
Gauss law (SCOPE \S3). This report answers each at classically reachable sizes and states exactly what remains
hardware-gated (\S\ref{sec:hw}).

\paragraph{Process.} The work was built by Claude Code from briefs written by a multi-agent overseer (Kiro), and
each milestone was independently reviewed (\code{review/01}--\code{review/12}). Claude Code took over the
overseer role from Kiro on 8 October 2026; this report was self-reviewed, with an independent read-only sub-agent
re-deriving the load-bearing numbers in fresh code (\code{review/13}). The full test suite passes (231 tests at
the time of writing). Nothing has been pushed or run on hardware.

% =====================================================================================
\section{Model, conventions and methods}\label{sec:methods}

\subsection{Hamiltonian and conventions}
\begin{equation}
  H \;=\; -J\sum_{\langle ij\rangle} Z_iZ_j \;-\; h\sum_i X_i ,
\end{equation}
on periodic tori with row-major site index $n=y L_x+x$ (qubit 0 the most significant bit).
\begin{itemize}
\item \textbf{Square reference circuits} \tg{established, QASM-verified}. The published circuits apply, per step, an
  $R_x$ half-layer, $R_{zz}$ on every bond, and an $R_x$ half-layer, with $|R_{zz}|=|R_x|=0.5$ and a single
  $R_y(\theta)$ preparation, $\theta=\pi/6$ (quench 1) or $-10^\circ$ (quench 2). A literal gate-by-gate replay
  equals our evolution at $(J,h)=(+1,+2)$, $dt=0.25$ to $10^{-12}$ \tg{derived, tested}. The step-1 QASM has 112
  distinct $R_{zz}$ pairs, equal to the bond set of a periodic $7\times8$ square torus under the row-major map
  and not under the column-major one \mea.
\item \textbf{Splitting.} The circuit order (X half, ZZ, X half) differs from the core integrator's ZZ-outside
  Strang splitting at this coarse step: on $4\times4$, step~1, $Z_2=0.78$ vs $0.63$ \mea. All comparisons with the
  published circuits use the circuit order; the reference is the circuit H2 ran, not continuous-time $H$
  (Trotter-vs-exact gap up to 0.066 for quench 1 and 0.298 for quench 2 on $4\times4$ \mea).
\item \textbf{Frustrated model.} 2503.20870 writes $H=J\sum ZZ+h\sum X$ with antiferromagnetic $J>0$; in our
  sign the frustrated model is $J=-1$ \est. The triangular torus uses bonds along $(1,0),(0,1),(1,1)$
  (``square + one diagonal''), $h/|J|=0.3$, $dt\,|J|=0.25$.

  \textbf{Scope of the triangular work (decided 2026-10-15).} This is a deliberately chosen study of
  the \emph{weak-field} frustrated regime, $h\le0.5$. It does \textbf{not} reproduce the reference's
  triangular quench, whose released circuits are at $h/|J|=2$ --- $r_{zz}(0.5)$, $r_x(0.5)/r_x(1.0)$
  in every file, identical to the square quenches. The field is chosen on physics grounds, not from the
  reference: the sector-restricted quantum-dimer description is an effective theory to \emph{first order}
  in $h/|J|$, leakage out of the defect-free sector scales as about $h^2$ (measured exponent 2.08), and
  this project's own field scan finds the in-sector verdicts stable for $h\le0.3$ and straining at
  $h=0.5$ (26\% of the weight out of sector). At $h/|J|=2$ none of that holds: noiselessly the in-sector
  probability collapses to 0.13--0.63 within a \emph{single} Trotter step and the in-sector weight falls
  with system size (0.55 at $N=9$, 0.11 at $N=12$) \mea, so the constraint whose usefulness this thread
  tests barely exists there. Independently of the field, the released quench-4 circuits prepare $r_y(0)$
  on all 54 qubits with no $x$ gates, so they fix neither candidate winding sector. See
  \texttt{DECISIONS.md} DECIDED-1.
\item \textbf{Observables.} $Z_2=\langle(\sum_iZ_i)^2\rangle/N^2$ (second moment, not the connected variance);
  connected $\langle Z_iZ_j\rangle_c$ binned by distance; dimer occupations $n_b=(1+Z_iZ_j)/2$, triangle
  defects $D_t$, winding numbers $(W_x,W_y)$ and flippable spins $\Pi_s$ (\S\ref{sec:gauss}).
\end{itemize}

\subsection{Codes, backends and ceilings}
Table~\ref{tab:codes} lists the four codes. Logical TFIM steps are authored once as pytket segments and run
three ways: an in-repo state-vector interpreter (\emph{local-sim}, cross-checked against pytket's
\code{get_statevector} to $10^{-12}$), \emph{Selene} (Quantinuum's emulator stack, via generated Guppy
programs with real mid-circuit measurement and feed-forward discard), and \emph{Stim} as the fault-tolerance
verifier only. A \code{nexus} backend exists as a stub that raises by design; nothing contacts Quantinuum
without \code{allow_hardware=True} and explicit per-run approval with a recorded cost.
Size ceilings, fixed from memory probes, are 24 spins for state vectors and 20 for exact diagonalization (ED),
counting the full physical register; Stim and the configuration-space Monte Carlo are unbounded.

\begin{table}[h]\centering\small
\caption{The four codes as built. Physical qubits = data + 1 syndrome ancilla + 1 flag.}\label{tab:codes}
\begin{tabularx}{\linewidth}{lllL}
\toprule
code & parameters & logical rotation & role \\
\midrule
Iceberg~\cite{iceberg} & $[[k{+}2,k,2]]$ & weight-2 native \code{ZZPhase} & implemented path; recommended \\
Concatenated Iceberg & $[[16,4,4]]$ & weight-4 parity gadget & implemented path; $d=4$ \der{}+\mea \\
5-qubit & $[[5,1,3]]$ & weights 3 (field), 6 (bond) & built, dominated \\
Steane & $[[7,1,3]]$ & weights 3 (field), 6 (bond) & built, dominated \\
\bottomrule
\end{tabularx}
\end{table}

\paragraph{Concatenated $[[16,4,4]]$.} Outer qubit $j$ of outer block $a$ is inner logical $a$ of inner block
$j$ (Knill's C4/C6 block pattern~\cite{knill}). Any undetected nontrivial error touches $\ge2$ inner blocks with
weight $\ge2$ each, so $d=4$ \der; the enumerated distance is 4 and the circuit detects all 16{,}248 nonidentity
data errors of weight $\le3$ \mea. With Chao--Reichardt single flags~\cite{chao}, no 1, 2 or 3 circuit faults
give a silent logical error (exact $\le3$-fault search), and Stim's search returns fault distance 4; without
flags it is 2 \mea. These fault distances refer to the syndrome-extraction (memory) circuit; the encoded Trotter
step still has single-fault residuals at its rotation gadgets (\S\ref{sec:q4}).

\subsection{Noise models (all assumed)}
\begin{itemize}
\item \textbf{Generic baseline} \asm: $p_2$ per two-qubit gate, $p_1=p_m=p_{\rm init}=p_2/10$, no idle term.
\item \textbf{H2 prior} \asm{}/\infr: $p_2=6\times10^{-4}$ (2503.20870's average native two-qubit infidelity,
  $6(1)\times10^{-4}$ \est), $p_1=p_m=p_{\rm init}=6\times10^{-5}$; a fraction $\eta=0.05$ (bracket 0.03--0.08)
  of each two-qubit error is an extra $Z\otimes Z$, read off a plot of the reference's cycle-benchmarking
  channel (Fig.~S3) \infr; coherent over-rotation $\sin^2\epsilon$ is unpublished and therefore scanned,
  $\{0,10^{-5},10^{-4},10^{-3}\}$ \asm.
\item \textbf{Gauss-law thread} \asm: independent single-qubit Pauli errors after every Trotter step
  (depolarizing, dephasing or bit-flip type), $p\in\{10^{-3},3\times10^{-3},10^{-2}\}$.
\end{itemize}
Not modelled anywhere: leakage and loss, crosstalk, correlated multi-qubit errors beyond the two-qubit channel,
SPAM asymmetry, and (except as a twirled stand-in) coherent errors.

% =====================================================================================
\section{Q1 --- Logical versus physical qubits}\label{sec:q1}

\subsection{The reference data and the comparison harness}
The released dataset (Zenodo record 18487089, MD5-verified) was loaded with shape checks against the published
step lists, and the geometry and gate angles were re-checked against the published QASM before any
comparison (run aborts on mismatch). Table~\ref{tab:q1} shows the first quench-1 steps against a reduced-lattice
simulation at matched dimensionless parameters. The check now compares the \emph{signed}
$(r_{zz}, r_x, r_y)$ triple against the four gauge-equivalent sign patterns and requires sign
uniformity across bonds; earlier drafts of this sentence claimed signs were checked while the code
compared magnitudes only, which would have accepted a non-gauge pattern that shifts $Z_2$ by more than
$0.05$. The released circuits are in fact sign-uniform throughout ($r_{zz}(0.5)$, $r_x(0.5)/r_x(1.0)$ in
all 80 files), so no published comparison was affected.

\begin{table}[h]\centering\small
\caption{Quench 1 ($\theta=30^\circ$), $Z_2$ pair average $(NZ_2-1)/(N-1)=\overline{\langle Z_iZ_j\rangle}_{i\neq j}$.
``H2'' columns are \emph{2503.20870's published 56-qubit physical-qubit data} (error bars: dataset bootstrap
\code{_err}; whether $1\sigma$ or a CI half-width is \unv). ``Sim'' is our noiseless $4\times5$ circuit
simulation \mea. Source: \code{data/q1_compare.txt}.}\label{tab:q1}
\begin{tabular}{rrccc}
\toprule
step & $t|J|$ & sim $4\times5$ & H2 \code{zne+znr} (published) & H2 raw (published) \\
\midrule
1 & 0.25 & 0.761 & $0.732\pm0.009$ & $0.726\pm0.007$ \\
2 & 0.50 & 0.730 & $0.691\pm0.010$ & $0.678\pm0.008$ \\
4 & 1.00 & 0.695 & $0.623\pm0.017$ & $0.582\pm0.013$ \\
20 & 5.00 & 0.690 & $0.554\pm0.040$ & $0.379\pm0.016$ \\
\bottomrule
\end{tabular}
\end{table}

\textbf{What this does and does not show.} It establishes that the loader, conventions, geometry and parameters
reproduce the published circuits. It is a \emph{finite-size trend} comparison and nothing more: the
reduced-lattice spread $|Z_2(4\times4)-Z_2(4\times5)|$ is 0.061 (quench 1) and 0.331 (quench 2) \mea, so it
cannot separate residual device noise from finite-size effects, and it contains \textbf{no logical-qubit
number}. A same-system check is possible at step 1 (light cone of one bond) but was not built.

\subsection{The encoded machinery}
The Iceberg-encoded TFIM step maps logical $e^{-i\phi\bar Z_i\bar Z_j}$ to the native \code{ZZPhase} on
$(i,j)$ and logical $e^{-i\phi\bar X_i}$ to $e^{-i\phi X_tX_i}$, authored as $H\otimes H$, \code{ZZPhase}, $H\otimes H$ because Selene's compiler rejects pytket's \code{XXPhase} \mea. A flagged $S_Z$ and $S_X$ syndrome
round follows every step, with post-selection on all detectors.
\begin{itemize}
\item \textbf{Transparency:} at zero noise, encoded logical $Z_2$ equals the bare simulation to $\le6\times10^{-15}$
  (steps 1--6, $2\times3$) \mea; for all four codes, Stim and local-sim agree on the \emph{full} logical state (all
  $4^k$ logical Pauli expectations) to $\le6\times10^{-14}$ at a Clifford instance \mea.
\item \textbf{Generic-noise execution} \cp: on $2\times3$ ($[[8,6,2]]$, 10 physical qubits), by step 6 bare $Z_2$
  drifts $\approx4.5\sigma$ ($p_2=10^{-3}$) and $\approx8\sigma$ ($3\times10^{-3}$) from ideal while post-selected encoded stays
  within $\approx2\sigma$, at 1.3--3$\times$ the shots; at $p_2=10^{-2}$ only $\approx5\%$ of shots survive and no
  advantage is resolved \mea. The advantage is not uniform: in 7 of the 24 $(p_2,\text{step})$ cells at
  $p_2\le3\times10^{-3}$ encoded sits further from ideal than bare, including steps 4 and 5, though
  encoded minus bare is within $1.0\sigma$ in every one of them; encoded beats bare at $\ge2.5\sigma$ only
  at step 6. (The $\approx2\sigma$ above is encoded against \emph{ideal}; one cell reaches $2.8\sigma$.)
\end{itemize}

\subsection{Does encoding help at all? The \texorpdfstring{$\eta$}{eta} crossover (FLAG-0)}\label{sec:flag0}
A two-qubit fault after a logical $e^{-i\phi P_aP_b}$ on the support $(a,b)$ that is $X_aX_b$, $Y_aY_b$ or $Z_aZ_b$
commutes with both Iceberg stabilizers and is a nontrivial logical: 3 of the 15 two-qubit Paulis, i.e.\ 1/5 of a
uniform gate fault, is undetectable \der{}+\mea{} (on the $2\times3$ step Stim finds 63 distinct undetected
Paulis, all $XX$/$YY$/$ZZ$ on the supports of the 21 logical gates; \code{docs/encoded_tfim.md} \S2 and its test.
This is a count of distinct Pauli \emph{operators} and is not the same quantity as the ``undetected-logical''
column of \code{data/encoded_iceberg.txt}, which counts detector-error-model fault \emph{mechanisms}). If a fraction $\eta$ of the error is extra $Z\otimes Z$, the undetectable rate per
logical gate is $p_2[(1-\eta)\tfrac{3}{15}+\eta]=p_2(0.2+0.8\eta)$. On the $2\times2$ test lattice the Iceberg
step has 12 logical rotations and the bare step 4 two-qubit and 8 one-qubit gates, so the crossover
$12\,p_2(0.2+0.8\eta)=4p_2+8p_1$ with $p_1=p_2/10$ gives $\eta=0.25$ exactly for Iceberg \der{}
($\approx0.27$ for the $d\ge3$ codes, whose bias mapping is frame-approximate \infr).

\textbf{The crossover depends on the lattice} \der{} (found in the self-review, \code{review/13}). On a torus with
$B$ bonds and $N$ sites the Iceberg step has $B+2N$ logical rotations and the bare step $B$ two-qubit plus $2N$
one-qubit gates, so the crossover is $(B+0.2N)/(B+2N)=0.2+0.8\eta$. The $2\times2$ torus is degenerate ($B=N$) and the
most pessimistic: the crossover is $\eta\approx0.36$ on $2\times3$, $\approx0.44$ on large square tori ($B=2N$) and
$\approx0.55$ on triangular tori ($B=3N$), and the ratio at $\eta=0.05$ falls from 0.60 ($2\times2$) to 0.49, 0.44 and
0.375. The triangular value equals the repo's own Stim accounting (\code{data/q3_size.txt}: $0.0086/0.0230=0.37$) \mea.

\begin{table}[h]\centering\small
\caption{Undetected-logical probability per Trotter step relative to the bare circuit's unprotected error
($2\times2$, Stim first order) at the H2 prior $p_2=6\times10^{-4}$ \cp. Source: \code{data/noise_sensitivity.txt} \S1.}
\label{tab:prior}
\begin{tabular}{lcccc}
\toprule
scenario & Iceberg & Steane & 5-qubit & concat $d=4$\\
\midrule
uniform ($\eta=0$) & $\times0.50$ & $\times0.42$ & $\times0.42$ & $\times0.42$\\
\textbf{prior $\eta=0.05$}, $\sin^2\epsilon=0$ & $\mathbf{\times0.60}$ & $\times0.53$ & $\times0.52$ & $\times0.53$\\
bracket $\eta=0.03$ / $0.08$ & $\times0.56$ / $\times0.66$ & $\times0.48$ / $\times0.59$ & $\times0.48$ / $\times0.59$ & $\times0.48$ / $\times0.59$\\
$\eta=0.05$, $\sin^2\epsilon=10^{-4}$ & $\times0.72$ & $\times0.66$ & $\times0.66$ & $\times0.66$\\
$\eta=0.05$, $\sin^2\epsilon=10^{-3}$ & $\times0.92$ & $\times0.91$ & $\times0.91$ & $\times0.91$\\
\emph{stress test} $\eta=0.5$ (${\approx}10\times$ prior) & $\times1.50$ & $\times1.50$ & $\times1.50$ & $\times1.50$\\
\midrule
acceptance after 3 steps (prior) & 0.92 & 0.57 & 0.64 & 0.64\\
\bottomrule
\end{tabular}
\end{table}

\textbf{The prior is conservative.} In the model $\eta$ is an \emph{extra} $Z\otimes Z$ fraction on top of the
uniform 1/15 share, so $\eta=0.05$ means a total ZZ share of 0.11, about twice the read-off (0.04--0.05); and the
read-off channel is dominated by $IZ{+}ZI$, putting Iceberg's undetectable share ($XX{+}YY{+}ZZ$) near 0.05--0.07
instead of the model's 0.24 \infr. Both effects favour encoding; a run with the read-off channel is not built (T-10).

\textbf{Predicted conditional GO} \cp. At the prior, every code leaves less undetected error per step than bare
(Iceberg $\times0.60$; $\times0.56$--$0.66$ over the $\eta$ bracket). The twirled over-rotation stand-in adds the
same $12\sin^2\epsilon$ to code and bare, so it erodes the ratio toward 1 but cannot cross it \der. The real flip
risk is un-twirled coherent $\epsilon$, which adds in amplitude and is ZZ-aligned on a logical-ZZ gate; the model
cannot represent it. The project therefore decided that encoded H2 runs use randomized compiling, as the
reference experiment does (decision D-RC). The $\eta$ value rests on a visual read-off of one published figure
\infr, and $\epsilon$ is unpublished. \textbf{Only the H2 emulator or device can settle this} (\S\ref{sec:hw}).

\paragraph{Q1 status.} Encoded machinery, harness and prior-based prediction: done and reviewed. The matched
head-to-head comparison against the released data, with error bars: \pend{} --- hardware-gated.

% =====================================================================================
\section{Q4 --- Which code, and why}\label{sec:q4}

\paragraph{Fault-tolerance accounting.} For all four codes, single faults in the syndrome-extraction gadgets
produce \emph{zero} undetected logical errors \mea. The whole undetectable residual sits at the logical
rotations and is independent of distance \der{}+\mea: Iceberg's native gate leaves $3/15\cdot p_2$; every parity
gadget leaves its own generator through the rotation point, $2/15\cdot p_2+1/3\cdot p_1$, whether $d=2$, 3 or 4.
The honest statement is: \emph{the syndrome gadget is single-fault tolerant; the encoded Trotter step is not.}

\begin{table}[h]\centering\small
\caption{Budget against the $\approx55$-qubit H2 scale (physical qubits incl.\ ancilla and flag) and per-step
cost \mea. L = runs locally and fits H2; H = fits H2 only; x = over budget.}\label{tab:budget}
\begin{tabular}{lcccccc|cc}
\toprule
code & $1\times2$ & $2\times2$ & $2\times3$ & $2\times4$ & $3\times4$ & $4\times4$ & 2q ops/step & lattice\\
\midrule
Iceberg & 6 L & 8 L & 10 L & 12 L & 16 L & 20 L & 28 & $2\times2$\\
Steane & 16 L & 30 H & 44 H & 58 x & 86 x & 114 x & 108 & $1\times2$\\
5-qubit & 12 L & 22 L & 32 H & 42 H & 62 x & 82 x & 84 & $1\times2$\\
concat $d=4$ & 18 L & 18 L & 34 H & 34 H & 50 H & 66 x & 184 & $2\times2$\\
bare & & & & & & & 4 & $2\times2$\\
\bottomrule
\end{tabular}
\end{table}

\begin{table}[h]\centering\small
\caption{Undetected-logical probability per step, $p_2=10^{-3}$ generic model, Stim first order \cp.}\label{tab:resid}
\begin{tabular}{lccccc}
\toprule
lattice & Iceberg & Steane & 5-qubit & concat & bare\\
\midrule
$2\times2$ & $2.40\times10^{-3}$ & $2.00\times10^{-3}$ & $2.00\times10^{-3}$ & $2.00\times10^{-3}$ & $4.80\times10^{-3}$\\
$2\times3$ & $4.20\times10^{-3}$ & $3.50\times10^{-3}$ & $3.50\times10^{-3}$ & $3.50\times10^{-3}$ & $1.02\times10^{-2}$\\
acceptance after 3 steps ($2\times3$) & 0.83 & 0.21 & 0.28 & 0.23 & ---\\
\bottomrule
\end{tabular}
\end{table}

\paragraph{Decision.} \textbf{Iceberg is the recommended encoded baseline.} Under the assumed noise \cp, the $d=3$
and $d=4$ codes buy about $1.2\times$ on the residual and pay in acceptance (2.9--3.9$\times$ lower after 3 steps on
$2\times3$ at $p_2=10^{-3}$, more at higher $p_2$, but only $\approx1.9$--$2.3\times$ at $2\times2$, so the lattice
belongs with the ratio; 1.4--1.6$\times$ at the H2 prior) and in two-qubit gates
(\textbf{6.6--8.6$\times$} per step at a matched $2\times2$ lattice: Iceberg 28, concatenated 184, 5-qubit 192,
Steane 240). Table~\ref{tab:budget}'s own rows use different lattices per code, and the "3--7$\times$" of
earlier drafts came from reading across them; it \emph{understated} the cost of the higher-distance codes and
so ran against this section's own conclusion. None was discarded: all four
were built, verified (transparency, Stim-vs-local-sim agreement) and executed on Selene with mid-circuit measurement
and feed-forward, and in each run the executed acceptance follows the first-order prediction from below (Iceberg
highest, concat lowest) \mea. The expectation in the brief that the 5-qubit code would be
discarded on gate cost was not supported by the measured counts. The ranking is robust across eight noise
variations (ZZ bias, idle noise, over-rotation, combined) \mea{} \cp; whether \emph{any} encoding helps is not
(\S\ref{sec:flag0}).

\paragraph{Selene vs Stim \cp.} Selene keeps fewer shots than Stim predicts because its depolarizing model charges
$p_1$ on every native single-qubit operation including $R_z$ ($724$ native vs $188$ pytket-level gates,
$\times3.85$) and $p_{\rm init}$ on allocations \mea. Matching these counts explains 60--80\% of the gap; a
1--3$\sigma$ residual remains \infr. This is a noise-definition difference, not a bug, and it widens Iceberg's lead.
Whether H2's native $R_z$ is noisy is open (HW-4).

\paragraph{Open lever.} The $d\ge3$ codes were run detect-and-discard. Active correction would raise their
acceptance and could change the ranking, but cannot remove the rotation-gadget residual (HW-2, not built).

% =====================================================================================
\section{Q2 --- Frustration: triangular versus square}\label{sec:q2}

All numbers in this section are exact ground states (ED; no sampling error, no noise) \mea.

\paragraph{Method.} $H$ commutes with $P=\prod_iX_i$. For $h>0$ the ground state is unique, positive and
$P$-even (Perron--Frobenius) \der. On the $4\times4$ square torus the $P$-odd partner lies below numerical
resolution at $h=0.3$, and an unrestricted Lanczos vector comes out as an arbitrary symmetry-broken mixture
($|\langle Z_i\rangle|$ from $\approx0.01$ to $\approx0.5$ over repeated runs) \mea. The ground state is therefore
computed exactly in the $P$-even sector; its energy matches unrestricted ED to within $2.6\times10^{-12}$ at every
point \mea. At $h=0$ both lattices are degenerate and nothing is interpreted; the triangular $4\times4$ ground state
is also quasi-degenerate at $h=0.1$ (gap $4.7\times10^{-5}$), though still unique. Distances are minimal-image
Euclidean in each lattice's embedding.

\paragraph{A correction to the earlier comparison (T-19).} A square torus with an odd periodic length is not
bipartite, so its antiferromagnet is itself frustrated: $h=0$ energy per bond is $-1$ on $4\times4$ but $-2/3$ on
$3\times4$ and $-0.8$ on $4\times5$ \mea. The earlier $3\times4$ ``unfrustrated'' square reference was therefore
not clean; its $-0.667$ vs $-0.333$ energy-per-bond contrast is partly odd-length frustration. Within the ED ceiling
the only clean square torus with both lengths $\ge3$ is $4\times4$, which is used as the matched pair.

\begin{table}[h]\centering\small
\caption{Connected correlations on the matched $4\times4$ pair, $J=-1$ \mea. Source:
\code{data/frustration_correlations.txt}.}\label{tab:q2}
\begin{tabular}{ccc|ccc}
\toprule
$h$ & square $C(1)$ & square $C(2.83)$ & triangular $C(1)$ & triangular $C(\sqrt3)$ & triangular $C(2)$\\
\midrule
0.1 & $-0.999$ & $+0.999$ & $-0.332$ & $-0.055$ & $+0.442$\\
0.3 & $-0.994$ & $+0.994$ & $-0.325$ & $-0.045$ & $+0.413$\\
0.6 & $-0.977$ & $+0.977$ & $-0.268$ & $+0.029$ & $+0.196$\\
1.0 & $-0.936$ & $+0.936$ & $-0.204$ & $+0.053$ & $+0.068$\\
\bottomrule
\end{tabular}
\end{table}

\paragraph{Results.}
\begin{itemize}
\item \textbf{Nearest-neighbour suppression (the frustration signature).} Triangular $C(1)$ is 3--5$\times$ weaker
  than on the bipartite square, close to $-1/3$ (``one frustrated bond per triangle'') at small $h$. With
  $\langle Z\rangle=0$, $C(1)$ is exactly the ZZ part of the energy per bond, so this restates the energy
  signature in correlation form \der.
\item \textbf{Distance resolution.} Square correlations do not decay across the torus (finite-size N\'eel order);
  triangular long-distance correlations are weaker than the square's and fall quickly with $h$ (at small $h$
the triangular $C(2)$ on $4\times4$ exceeds its own $|C(1)|$). Equal $h$ is not equal distance from
  each lattice's critical field (square $h_c\approx3.04|J|$ \unv{}, recalled QMC value; triangular lower,
  $\approx1.6|J|$ in the Isakov--Moessner diagram~\cite{isakov} \unv{}, recalled), so part of the field fragility is
  that offset.
\item \textbf{Torus-shape dependence.} At $h=0.3$ the triangular $\sqrt3$ shell is $+0.46$ on $3\times4$ but
  $-0.04$ on $4\times4$, and the $r=2$ shell is $-0.69$ vs $+0.41$. The triangular embedding sits at the $J'=J$
  point of Iglovikov \emph{et al.}~\cite{iglovikov}, near a disorder line and incommensurate floating phase \est;
  no ``frustration monotonically weakens order'' story is told.
\end{itemize}

\paragraph{Not claimed.} No thermodynamic-limit or phase statement ($N=12$--16); \textbf{no clock-order claim} (the
three-sublattice order parameter, K3, was not built; the positive $\sqrt3$ shell on $3\times4$ is consistent with
clock correlations~\cite{isakov,humeniuk} but $C(r)$ cannot separate them from shape effects). The planned size
trend (Part B) was cancelled; within ED $\le20$ it would have had only one clean square point.

\paragraph{Frustration-specific observables.} The dimer, defect and winding observables of \S\ref{sec:gauss} exist
only on the triangular lattice; the defect density per triangle grows as $(h/|J|)^2$ (fitted slopes 1.96--1.98 in
the ground state) \mea.

% =====================================================================================
\section{Q3 --- Thermalization}\label{sec:q3}

\subsection{Per-sector ergodic baseline}
ETH in quantum dimer models holds only within a winding sector~\cite{lanpowell} \est, so the ergodic reference must
be per sector: on $3\times4$, $W_0=(0,2)$, the flippable density is 0.2667 in-sector vs 0.2353 globally \mea.
For $|h|\ll|J|$ the quench is generated at first order by $H_{\rm QDM}=-h\sum_s\Pi_sX_s$ inside the defect-free
manifold; its spectrum is symmetric about 0, so a product state sits at infinite temperature of the constrained
band \der. The baseline is therefore the \textbf{uniform ensemble over the flip-connected component of the initial
configuration $c_0$ in its defect-free sector}. On every torus enumerated ($N\le20$) each non-frozen sector is a
single component \mea; above $N=20$ this is \unv.

The Monte Carlo sampler flips a uniformly chosen site iff it is flippable. Each move is its own inverse with the same
probability, so the kernel is symmetric and uniform on the component \der; a $\chi^2$ test confirms uniformity on
$3\times3$, and the sector is re-checked every sweep. Errors use 32 independent chains (autocorrelation-aware; the
naive error would be 2--3$\times$ too small). MCMC agrees with brute-force enumeration within error on $3\times3$,
$3\times4$, $3\times6$ and $4\times5$ (max $|z|=1.39$ on discriminating observables) \mea, and runs at the
reference $9\times6$ size in seconds (flippable density $0.3026\pm0.0002$).

\subsection{Ideal relaxation and the finite-size offset}
\begin{table}[h]\centering\small
\caption{$3\times6$, three-sublattice start, $W_0=(0,0)$: late-window sector-conditioned values (noiseless
circuit) vs the infinite-temperature sector baseline $B_{W_0}$ and the exact $H_{\rm QDM}$ diagonal ensemble
(DE) \mea. $\pm$ is the temporal spread over the window. Source: \code{data/q3_relaxation.txt}.}\label{tab:q3}
\begin{tabular}{lcccccc}
\toprule
observable & $B_{W_0}$ & DE & $h=0.1$ & $h=0.2$ & $h=0.3$ & $h=0.5$\\
\midrule
flippable density & 0.3310 & 0.3880 & $0.383\pm0.018$ & $0.376\pm0.017$ & $0.384\pm0.019$ & $0.383\pm0.018$\\
$C(0.866)$ & $+0.110$ & $+0.128$ & $+0.127$ & $+0.125$ & $+0.128$ & $+0.127$\\
late $P_{\rm out}$ & & & $6.2\times10^{-3}$ & $2.5\times10^{-2}$ & $5.9\times10^{-2}$ & $\mathbf{0.26}$\\
\bottomrule
\end{tabular}
\end{table}

The curves relax (the start has flippable density 0.667), but to the finite-cluster diagonal ensemble, not to the
infinite-temperature sector value. The dynamics conserves $\langle H_{\rm QDM}^2\rangle=h^2\times$(number of
flippable sites), and the three-sublattice start carries $2\times$ the uniform ensemble's energy variance \der.
The offset shrinks from $3\times3$ to $3\times6$ (0.092 $\to$ 0.057) \mea{} and is expected to vanish as
$N\to\infty$ \infr, not demonstrated beyond $N=18$. Leakage out of the sector is $O((h/|J|)^2)$; the exact bound
$|\langle O\rangle-\langle O\rangle_{\rm in}|\le P_{\rm out}\,{\rm range}(O)$ holds at every step \der. At $h=0.5$ a
quarter of the weight has left the sector: a stress point, not a result.

\begin{figure}[h]\centering
\maybefig{q3_relaxation.png}{0.9\linewidth}
\caption{Ideal (noiseless, unencoded) relaxation on $3\times6$ from the three-sublattice state at $h=0.1$ and 0.3,
raw and sector-conditioned, against $B_{W_0}$ and the DE, with the exact leakage envelope. Classical simulation;
not an H2 result.}\label{fig:relax}
\end{figure}

\subsection{Is \texorpdfstring{$dt=0.25$}{dt = 0.25} below the Trotter threshold?}
Local observables stay faithful below a sharp Trotter-step threshold~\cite{heyl}. Here the threshold is the
Floquet resonance $dt=\pi/(|J|n_s)$, where flipping a non-flippable spin costs $2|J|n_s$ ($n_s\in\{2,4,6\}$);
$\pi/6=0.524$ for the $n_s=6$ spins of the starts \der, broadened downward by $h$. Measured thresholds (absolute criterion: time-max flippable
error $0.05$) are 0.38--0.52 at $h\le0.3$ and 0.32--0.42 at $h=0.5$ \mea. Second order is confirmed (local slopes
1.97--2.04, infidelity $\propto dt^4$). \textbf{For $h\le0.3$, $dt=0.25$ sits 1.3--2.0$\times$ below the threshold}
(relative and absolute criteria); at $h=0.5$ it is past the relative threshold. The threshold location is
invariant under the choice of Strang ordering, since the two step operators are unitarily conjugate and so share
their Floquet quasienergies; re-measured in the published circuits' ordering, every bracket is equal or higher,
so this margin is conservative. The sector-conditioned verdict quantities have Trotter bias $\le0.12$ of their
tolerance in either ordering.

The leakage and raw defect density, however, carry a Trotter bias \textbf{whose sign depends on the ordering}.
In this project's splitting (\code{src/trotter.py}, $ZZ$ outside) $P_{\rm out}$ is inflated $\approx10$--25\% and
the $3\times6$ defect density $\approx40\%$ (a blend of Floquet factors, predicted 1.39--1.44$\times$ vs measured
1.39--1.43$\times$) \der{}+\mea. In the \emph{published circuits'} ordering ($X$ half, $ZZ$, $X$ half, as in
\code{src/compare.py}) the same quantities are \emph{deflated}: $P_{\rm out}$ by 19--35\% and the defect density
by 25--51\% ($3\times6$ full-window ratio 0.605 against 1.389) \mea. Earlier drafts stated that the absolute
defect densities at $dt=0.25$ ``describe the circuit, not continuous time'': that had the sign backwards. They
describe this project's own splitting. Any statement about the reference circuit must use the circuit ordering;
see \code{data/trotter_order_compare.txt}.

\begin{figure}[h]\centering
\maybefig{trotter_threshold.png}{0.85\linewidth}
\caption{Time-maximal Trotter error vs $dt$ with the $dt^2$ guide, the $\pi/6$ and $\pi/4$ resonances, the
threshold criterion and the working $dt=0.25$. Classical, noiseless.}\label{fig:trotter}
\end{figure}

\subsection{Noisy and Iceberg-encoded curves \texorpdfstring{\cp}{(classical prediction, assumed noise)}}
Under the H2-prior noise model, both arms run the same circuit-order Strang circuit; the readout distribution is
computed exactly at every step from Pauli-fault trajectories (errors: batch-means over fault patterns).

\begin{table}[h]\centering\small
\caption{$3\times4$, flippable start, $h=0.3$: mean $|O_{\rm noisy}-O_{\rm ideal}|$ of the flippable density, and the
shot-cost factor $1/\text{acceptance}$ \cp. Gauss modes are scored against the ideal sector-conditioned curve.
Source: \code{data/q3_noisy.txt}.}\label{tab:noisy}
\begin{tabular}{lcc}
\toprule
mode & early ($t\le12$) & late ($t=20$--40)\\
\midrule
bare, raw & $0.052\pm0.005$ ($\times1$) & $0.173\pm0.005$ ($\times1$)\\
bare + Gauss-law post-selection & $0.0017\pm0.0004$ ($\times1.43$) & $0.0083\pm0.0011$ ($\times4.8$)\\
Iceberg, post-selected & $0.019\pm0.008$ ($\times3.1$) & not reachable\\
Iceberg + Gauss & $0.0009\pm0.0009$ ($\times3.5$), statistical floor & not reachable\\
\bottomrule
\end{tabular}
\end{table}

\begin{itemize}
\item Unmitigated, the bare curve heats toward the full-space infinite-temperature value ($\approx3.7$ faults per
  shot over the window).
\item \textbf{Gauss-law post-selection is the cheap, effective mitigation} for these diagonal observables:
  $\approx30\times$ early and $\approx21\times$ late on $3\times4$ against the conditioned target (5.5--6$\times$ on
  $3\times3$), for $\times1.2$--4.8 shots.
\item Iceberg alone gives $\times2.7$ (1$\sigma$ range $\approx1.9$--4.7) for $\times3.1$ shots (matching the Stim per-step rate ratio) and its survival
  (8\% by $t=12$) confines it to early times; Iceberg + Gauss is at the statistical floor, not resolved as better
  than Gauss alone. In this model the encoding's marginal value over the model's own Gauss law is at most modest.
\item \textbf{Noise can mimic thermalization (T-17).} The noisy conditioned late values drift from the ideal
  DE-like value toward $B_{W_0}$ ($3\times3$: 0.427 noisy vs 0.477 ideal, $B_{W_0}=0.381$). An H2 curve landing on
  $B_{W_0}$ is therefore not by itself evidence of thermalization; the noise must be part of the comparison.
\end{itemize}

\begin{figure}[h]\centering
\maybefig{q3_noisy.png}{0.85\linewidth}
\caption{$3\times4$ ideal vs bare, bare + Gauss, Iceberg and Iceberg + Gauss, with both acceptances. Noisy curves:
classical prediction under an assumed noise model (H2 prior); not an H2 result.}\label{fig:noisy}
\end{figure}

\subsection{Size story}
The sector baseline follows to $12\times12$ (flippable density $0.2947\pm0.0002$); a weighted $a+b/N$ fit on
$N\ge36$ gives $a\approx0.290$ \asm{} fit form, \infr{} extrapolation, so the $9\times6$ reference value still sits
$\approx0.013$ above it. The noise budget is linear in $N$ \der: per step the bare fault probability is
$1.92\times10^{-3}N$, i.e.\ $\approx17$ faults per shot over the 160-step window at $9\times6$, and Iceberg
survival reaches 1\% by $t\approx5$. \textbf{Iceberg on $9\times6$ needs 58 physical qubits, more than H2's 56}
(T-18), so an encoded Q3 run must use a smaller torus (e.g.\ $6\times6$, 40 physical qubits).

\begin{figure}[h]\centering
\maybefig{q3_size.png}{0.85\linewidth}
\caption{Sector baseline vs $1/N$ with the inferred fit, and noise-budget times vs $N$. Classical statistics and
first-order accounting under assumed noise.}\label{fig:size}
\end{figure}

\begin{figure}[h]\centering
\maybefig{q3_field_scan.png}{0.85\linewidth}
\caption{Late conditioned values vs $h$ (with $B_{W_0}$ and DE) and late $P_{\rm out}$ vs $h$. Classical, noiseless.}
\label{fig:field}
\end{figure}

\paragraph{Not claimed.} No thermodynamic-limit thermalization statement: the quantum curves stop at $N=18$, where
the finite-size offset dominates.

% =====================================================================================
\section{The emergent Gauss law as a mitigation resource (SCOPE \S3)}\label{sec:gauss}

\paragraph{Structure \der.} With $n_b=(1+Z_iZ_j)/2$, each triangle has $G_t=\sum_{b\in t}n_b\in\{1,3\}$; the defect
$D_t=(G_t-1)/2$; the Ising energy is $E=|J|(2\sum_tD_t-N)$, so the classical ground manifold is exactly the
defect-free (Gauss-law) subspace, the Moessner--Sondhi dimer picture~\cite{moessnersondhi}. Defects come in
pairs, winding numbers are sector labels on the defect-free subspace, and $H_{\rm QDM}$ conserves both exactly;
$H_{\rm TFIM}$ violates them at $O((h/|J|)^2)$ (fitted slopes 1.96--2.05; stripe prefactor within 1.1\% of first
order) \mea. The emergent constraint is a local $\mathbb Z_2$ Gauss law~\cite{msf}.

\paragraph{Detection power.} $Z$ errors commute with every $D_t$, $n_b$ and $W$: invisible when they occur, and the
defect response to dephasing is $\approx70\times$ weaker than to $X$/$Y$ \mea. $X$/$Y$ errors are detected with
probability 0.46--0.54; a flip of a flippable spin is a legal move and is invisible. As a code it has distance 1.

\begin{table}[h]\centering\small
\caption{Verdict (noise model assumed, per-qubit Pauli after each step; $3\times3$ exact density matrix, $3\times4$
trajectories). Source: \code{docs/gauss_law_dimers.md} \S5, \code{data/gauss_law_mitigation.txt}.}\label{tab:gauss}
\begin{tabularx}{\linewidth}{P{4.2cm}P{3.3cm}L}
\toprule
use & verdict & basis\\
\midrule
error correction & \textbf{no} (passive constraint) & distance 1; no recovery map; ideal dynamics violates the constraint. Engineered monitored $\mathbb Z_2$ schemes with feedback exist~\cite{wauters} but were not modelled. \\
post-selection on defects & \textbf{modest yes}, $X$/$Y$-type noise only & 2--8$\times$ bias reduction on spin correlators at $p\le3\times10^{-3}$ for 1.1--1.8$\times$ shots; floored by an $O((h/J)^2)$ physics bias ($\approx0.013$ at $h/|J|=0.3$); acceptance falls with size even noiseless \mea\\
\quad under dephasing & \textbf{no} & post-selection only flips the residual's sign \mea\\
self-mitigation by depolarizing rescaling & \textbf{no} as a general method & under local noise the defect-derived fidelity is not the per-observable decay; up to $32\times$ worse on some observables (the maximum in \code{data/gauss_law_mitigation.txt}: bit-flip $p=10^{-2}$ at $t=3.25$; $25\times$ for depolarizing $p=3\times10^{-3}$) \mea\\
exact-symmetry methods (symmetry expansion, GUESS, symmetry-adjusted shadows) & \textbf{not applicable} & a $Z$-product quench has $\langle\prod X\rangle=0$, and Pauli channels keep it 0; TFIM lacks the needed U(1) \der\\
platforms & \textbf{undetermined} & depends on each platform's residual $X$/$Y$ fraction; Gauss-law post-selection has been realized on H2 for an \emph{engineered} one-form syndrome~\cite{sukeno} \tg{established, abstract}\\
\bottomrule
\end{tabularx}
\end{table}

\textbf{One line:} mostly \emph{no}; a narrow, conditional \emph{yes} for post-selection against bit-flip-type
residual noise, floored at $O((h/J)^2)$. ZNE is the practical baseline (the reference used it); PEC costs
$\gamma^2\approx5.7\times10^3$ for 54 qubits over 40 steps (the thread's $t=10$ window) at $p=10^{-3}$ under this
model \mea.
In Q3 (\S\ref{sec:q3}) the same post-selection is the effective mitigation for in-sector diagonal observables.

% =====================================================================================
\section{Findings versus questions}\label{sec:table}

\begingroup\small
\begin{longtable}{P{1.2cm}P{5.2cm}P{4.6cm}P{3.6cm}}
\caption{Findings vs questions (SCOPE \S13). ``Done'' means evidence with uncertainties at classically reachable
sizes; hardware-gated items are marked. Evidence paths are relative to the repository root.}\label{tab:fvq}\\
\toprule
Q & finding & evidence & status\\
\midrule\endfirsthead
\toprule
Q & finding & evidence & status\\
\midrule\endhead
\bottomrule\endfoot
Q1 & Conventions, geometry and parameters reproduce 2503.20870's published circuits; reduced-lattice trend vs released data, no logical-vs-physical number & \code{src/compare.py}, \code{data/q1_compare.txt}, \code{docs/q1_compare.md} & \mea{} --- done (harness)\\
Q1 & Encoded TFIM is transparent at zero noise; under generic assumed noise, encoded tracks ideal longer than bare at $p_2\le3\times10^{-3}$ on $2\times3$, at 1.3--3$\times$ shots & \code{data/encoded_iceberg.txt}, \code{docs/encoded_tfim.md} & \mea, \cp\\
Q1 & ``Encoding helps'' flips at ZZ fraction $\eta\approx0.25$ on the $2\times2$ test lattice ($\approx0.36$--0.55 on larger tori); at the conservative H2 prior $\eta\approx0.05$, Iceberg $\times0.60$ of bare on $2\times2$: predicted conditional GO, given randomized compiling & \code{data/noise_sensitivity.txt}, \code{research/06-h2-noise-prior.md} & \der{}+\mea, \cp; prior \infr\\
Q1 & Matched logical-vs-physical comparison against the released data, with error bars & --- & \pend{} \textbf{hardware-gated} (HW-5, G1--G3)\\
\midrule
Q4 & All four codes built, verified, executed; residual set by rotation gadgets, not distance & \code{src/encoded_codes.py}, \code{data/encoded_codes.txt} & \der{}+\mea\\
Q4 & Iceberg recommended; $d\ge3$ buy $\approx1.2\times$ residual for 3--4$\times$ lower acceptance ($p_2=10^{-3}$; 1.4--1.6$\times$ at the H2 prior); ranking robust to noise variations & \code{docs/encoded_tfim.md} B5, \code{data/noise_sensitivity.txt} & \mea, \cp\\
Q4 & Concatenated $[[16,4,4]]$: $d=4$, syndrome-circuit fault distance 4 with flags (2 without) & \code{docs/concatenated_iceberg_d4.md}, \code{data/concat_d4.txt} & \der{}+\mea\\
Q4 & H2-calibrated ranking; active correction for $d\ge3$ & --- & \pend{} hardware-gated (G2, G4)\\
\midrule
Q2 & Matched $4\times4$: triangular $C(1)$ 3--5$\times$ weaker than bipartite square; square long-range correlation does not decay, triangular weaker and field-fragile & \code{src/frustration.py}, \code{data/frustration_correlations.txt}, \code{docs/q2_frustration.md} & \mea{} --- done at reachable sizes\\
Q2 & Odd-length square tori are frustrated; the old $3\times4$ reference is corrected (T-19) & same; \code{data/frustration_compare.txt} & \mea\\
Q2 & Triangular long-distance signs depend on torus shape (disorder-line caveat); no clock-order claim & same; \cite{iglovikov} & \mea; K3 not built\\
Q2 & Dimer/defect/winding observables, defect density $\propto(h/|J|)^2$ & \code{src/dimers.py}, \code{data/gauss_law_scaling.txt} & \der{}+\mea\\
\midrule
Q3 & Per-sector MCMC baseline, exact within error up to $N=20$, scalable to $12\times12$ & \code{src/mcmc_baseline.py}, \code{data/mcmc_baseline.txt} & \der{}+\mea\\
Q3 & Ideal curves relax to the finite-cluster DE, not the infinite-$T$ sector value (finite-size offset) & \code{data/q3_relaxation.txt}, \code{docs/q3_thermalization.md} & \mea; $N\to\infty$ \infr\\
Q3 & $dt=0.25$ is 1.3--2$\times$ below the Trotter threshold for $h\le0.3$; $h=0.5$ a stress point & \code{src/trotter_threshold.py}, \code{data/trotter_threshold.txt} & \der{}+\mea\\
Q3 & Gauss post-selection $\approx5.5$--30$\times$ for $\times1.2$--4.8 shots; Iceberg $\times2.7$ for $\times3.1$; noise can mimic thermalization & \code{src/noisy_q3.py}, \code{data/q3_noisy.txt} & \cp\\
Q3 & Iceberg $9\times6$ needs 58 $>$ 56 qubits; $\approx17$ faults/shot bare at $9\times6$ & \code{data/q3_size.txt} & \der{} accounting, \cp\\
Q3 & Matched H2 relaxation curves & --- & \pend{} hardware-gated\\
\midrule
Gauss & Correction: no. Detection: modest yes for $X$/$Y$ noise, floored at $O((h/J)^2)$. Dephasing: no. Self-mitigation: no. Exact-symmetry methods: not applicable. Platforms: undetermined & \code{src/gauss_mitigation.py}, \code{data/gauss_law_mitigation.txt}, \code{docs/gauss_law_dimers.md} & \der{}+\mea, assumed noise\\
\end{longtable}
\endgroup

% =====================================================================================
\section{What is hardware-gated, and the plan}\label{sec:hw}

Nothing below has been run. Each step needs \code{allow_hardware=True}, the owner's explicit per-run approval
and a recorded cost estimate. Every result will be labelled classical prediction, H2 \emph{emulator} or H2
\emph{hardware}, never mixed; Stim remains the fault-tolerance verifier.
\begin{enumerate}
\item \textbf{G1 / HW-5 --- calibration.} Iceberg $2\times2$, one Trotter step, on the H2 emulator; compare discard
  rate and logical $Z_2$ with the Selene and Stim predictions before any hardware cost.
\item \textbf{G2 --- the inputs of FLAG-0.} Obtain $\eta$ and $\epsilon$ for H2's native gate from the best
  labelled source and rerun the sensitivity table.
\item \textbf{G3 --- the decisive test.} Encoded (Iceberg, randomized-compiled) vs bare on H2 at matched shots,
  lattice, steps and observable; a yes / no / inconclusive verdict on ``encoding helps''. The lattice must be
  one where \emph{both} arms fit in H2's 56 qubits: $6\times8$ (52 physical) is the closest in shape to the
  reference's $7\times8$, which cannot be encoded at all ($k+4=60>56$). The released physical data enters as an
  external reference --- quantitatively at step 1 through an exact light-cone finite-size correction, and as a
  trend thereafter --- not as the bare arm. The same reduction was already made for Q3 and is now explicit
  for Q1 (\texttt{DECISIONS.md} OPEN-2).
\item \textbf{G4 --- active correction} for $d\ge3$, designed and verified in Stim and Selene first.
\end{enumerate}
Two decisions from this work apply to the hardware runs. An encoded Q3 run cannot use the $9\times6$ reference torus on
56 qubits (T-18), so it runs on $6\times6$ (40 physical qubits) and on the smallest relevant torus, $3\times4$
(16 physical; the smallest with even $k$ and more than one winding sector, with all classical predictions in hand).
A Q3 curve reaching $B_{W_0}$ is not called thermalization unless it is compared with the noise (T-17); at $6\times6$
no classical noisy prediction is possible ($36>24$ spins), so that run includes a noise-scaling arm. Under the assumed
prior the $6\times6$ Iceberg arm reaches only early times (survival 10\% at $t\approx3.8$) \cp.

% =====================================================================================
\section{Limitations}\label{sec:limits}
\begin{itemize}
\item All noisy results use assumed models. The H2 prior's $\eta$ is a visual read-off of one published figure; the
  coherent over-rotation is unpublished; the Selene and Stim noise definitions differ by native-gate counting.
\item Sizes: executed encoded runs use 2--6 logical qubits; exact Q2 ground states $N\le16$; quantum Q3 curves
  $N\le18$; larger sizes appear only in configuration-space statistics and first-order accounting.
\item The Q1 table is a reduced-lattice trend, not a same-system or logical-vs-physical comparison.
\item Open conventions in the reference: the dimer-correlation distance is a graph distance in 2503.20870 while the
  repo bins by Euclidean distance (affects only direct comparison with its $C_r$ curves); whether the dataset's
  error bars are $1\sigma$ is \unv.
\item Not built (owner's parked decisions): active gauge protection (K1), a clock-order observable (K3), a diffusion
  observable (K5); the Q2 size trend was cancelled.
\end{itemize}

% =====================================================================================
\appendix
\section{Reproducibility}
All settings live in \code{config.py}. Scripts are run as \code{.venv/bin/python -I scripts/<name>.py} and write
only to \code{data/} (figures to \code{figures/}). Environment: Python 3.12.14; numpy 2.5.3, scipy 1.18.1, stim 1.16.0,
pytket 2.18.4, guppylang 1.1.2, selene-sim 0.3.4 (pinned in \code{requirements.txt}). Test suite:
\code{python -m pytest tests/ -q}.

\begin{center}\small
\begin{tabularx}{\linewidth}{LL}
\toprule
data file (\code{data/}) & produced by (\code{scripts/})\\
\midrule
\code{q1_compare.txt} & \code{run_compare.py --data <Zenodo dir>}\\
\code{encoded_iceberg.txt}, \code{encoded_codes.txt} & \code{run_encoded.py}, \code{run_encoded_codes.py}\\
\code{noise_sensitivity.txt}, \code{selene_gap.txt} & \code{run_noise_sensitivity.py}, \code{diagnose_selene_gap.py}\\
\code{concat_d4.txt} & \code{run_concat_d4.py}\\
\code{frustration_compare.txt}, \code{frustration_correlations.txt} & \code{compare_frustration.py}, \code{compare_frustration_correlations.py}\\
\code{gauss_law_scaling.txt}, \code{gauss_law_mitigation.txt} & \code{run_gauss_law.py}\\
\code{mcmc_baseline.txt}, \code{q3_relaxation.txt} & \code{run_q3_baseline.py}\\
\code{trotter_threshold.txt} & \code{run_trotter_threshold.py}\\
\code{q3_noisy.txt}, \code{q3_size.txt} & \code{run_q3_noisy.py}, \code{run_q3_size.py}\\
\code{figures/*.png} & \code{plot_q3.py}\\
\bottomrule
\end{tabularx}
\end{center}

% =====================================================================================
\begin{thebibliography}{99}\small
\bibitem{haghshenas} R.~Haghshenas \emph{et al.}, ``Digital quantum magnetism on a trapped-ion quantum computer,''
  Nature \textbf{653}, 56--62 (2026); arXiv:2503.20870. Data: Zenodo record 18487089.
\bibitem{froland} Froland \emph{et al.}, ``Realizing error suppression in partially fault-tolerant quantum
  simulations with IBM quantum computers,'' arXiv:2607.24947 (2026).
\bibitem{helios} Quantinuum, ``Computing with many encoded logical qubits beyond break-even,'' arXiv:2602.22211.
\bibitem{begusic} T.~Begu\v{s}i\'c, J.~Gray, G.~K.-L.~Chan, Sci.\ Adv.\ \textbf{10}, eadk4321 (2024); arXiv:2308.05077.
\bibitem{tindall} J.~Tindall \emph{et al.}, PRX Quantum \textbf{5}, 010308 (2024); arXiv:2306.14887.
\bibitem{iceberg} C.~N.~Self, M.~Benedetti, D.~Amaro, ``Protecting expressive circuits with a quantum error
  detection code,'' arXiv:2211.06703.
\bibitem{knill} E.~Knill, ``Quantum computing with realistically noisy devices,'' Nature \textbf{434}, 39--44 (2005).
\bibitem{chao} R.~Chao, B.~W.~Reichardt, ``Quantum error correction with only two extra qubits,'' Phys.\ Rev.\ Lett.\
  \textbf{121}, 050502 (2018); arXiv:1705.02329.
\bibitem{moessnersondhi} R.~Moessner, S.~L.~Sondhi, ``Ising models of quantum frustration,'' Phys.\ Rev.\ B
  \textbf{63}, 224401 (2001).
\bibitem{msf} R.~Moessner, S.~L.~Sondhi, E.~Fradkin, Phys.\ Rev.\ B \textbf{65}, 024504 (2002); arXiv:cond-mat/0103396.
\bibitem{isakov} S.~V.~Isakov, R.~Moessner, Phys.\ Rev.\ B \textbf{68}, 104409 (2003); arXiv:cond-mat/0302105.
\bibitem{humeniuk} S.~Humeniuk, Phys.\ Rev.\ B \textbf{93}, 104412 (2016); arXiv:1602.02979.
\bibitem{iglovikov} V.~I.~Iglovikov \emph{et al.}, Phys.\ Rev.\ B \textbf{87}, 214415 (2013); arXiv:1303.0883.
\bibitem{lanpowell} Z.~Lan, S.~Powell, ``Eigenstate thermalization hypothesis in quantum dimer models,'' Phys.\ Rev.\
  B \textbf{96}, 115140 (2017); arXiv:1706.02601.
\bibitem{heyl} M.~Heyl, P.~Hauke, P.~Zoller, ``Quantum localization bounds Trotter errors in digital quantum
  simulation,'' Sci.\ Adv.\ \textbf{5}, eaau8342 (2019); arXiv:1806.11123.
\bibitem{wauters} Wauters \emph{et al.}, ``Symmetry-protection Zeno phase transition in monitored lattice gauge
  theories,'' arXiv:2405.18504.
\bibitem{sukeno} H.~Sukeno, E.~Rinaldi, S.~Okuda, measurement-based $\mathbb Z_2$ lattice gauge theory on H2,
  arXiv:2608.04290 (2026) [abstract read].
\end{thebibliography}

\end{document}
